% !TEX root=../main.tex
\section{Related Works}\label{sec:related-work}

\noindent
\textbf{Root Cause Analysis.}
% \paragraph{Root Cause Analysis}
Corollary~\ref{thm:l2-necessary} explains that graph construction is a common step in the RCA literature for online service system operation.
DFS-based methods~\cite{Chen:2014,Lin:2018,Liu:2021} traverse abnormal sub-graph, which is sensitive to anomaly detection results.
Some works adopt random walk~\cite{Wang:2018,Meng:2020,Ma:2020} or PageRank~\cite{Wang:2021} to score candidate root cause indicators, lacking explainability.
% Random walk-based and PageRank-based methods~\cite{Wang:2018,Meng:2020,Ma:2020,Wang:2021} lack explainability.
Another line of works is invariant network-based methods~\cite{Cheng:2016,Ni:2017}.
As these works adopt the pair-wise manner to learn the invariant relations, it is hard for them to reach the knowledge of RCA, restricted by CHT~\cite{Bareinboim:2022}.
No methods above utilize causal inference.
Sage~\cite{Gan:2021} conducts counterfactual analysis to locate root causes without a formal formulation.
Corollary~\ref{thm:l3-unnecessary} states that counterfactual analysis is unnecessary.
Hence, we did not include this method as a baseline.
Meanwhile, CHT~\cite{Bareinboim:2022} indicates that it can be hard to conduct counterfactual analysis even with a CBN.

The definition of root cause analysis varies with the scenario in the literature.
% We treat the observed projection of a fault as the desired answer.
% In contrast, some applications require an answer beyond the data, taking RCA as a classification task with supervised learning~\cite{Yi:2021}.
% For homogeneous devices or services, operators are interested in the common features~\cite{Zhang:2021}.
% Accordingly, multi-dimensional root cause analysis is conducted.
Some applications require an answer beyond the data, taking RCA as a classification task with supervised learning~\cite{Yi:2021}.
For homogeneous devices or services, operators are interested in the common features~\cite{Zhang:2021}.
Accordingly, a multi-dimensional root cause analysis is conducted.
In contrast, we treat the observed projection of a fault as the desired answer.
The \nameref{thm:rc-criterion} further relates RCA in this work with \textit{contextual anomaly detection}~\cite{Chandola:2009}, treating parents in the CBN as the context for each variable.
We take complex contextual anomaly detection methods as future work.

\noindent
\textbf{Causal Discovery.}
% \paragraph{Causal Discovery}
The task to obtain the CBN is named \textit{causal discovery}.
We refer the readers to a recent survey~\cite{Guo:2020} for a thorough discussion.
NOTEARS~\cite{Zheng:2018} converts the DAG search problem from the discrete space into a continuous one.
Following NOTEARS, some recent works are based on gradient descent~\cite{He:2021}.

Although causal discovery has its sound theory, the CBN discovered from data directly is not explainable for human operators.
In contrast, some works obtain the CBN based on domain knowledge.
MicroHECL~\cite{Liu:2021} traces the fault along with traffic, latency, or error rate in the call graph.
Meanwhile, Sage~\cite{Gan:2021} constructs the CBN among latency and machine metrics.
The structural graph proposed in this work is compatible with the assumptions in these two works, extending the kinds of meta metrics.
